% GENERATED FILE. Edit canonical lessons, then run npm run generate:books.
% canonical-source-hash: fnv1a64:d3f5678ddcaa461f
% canonical-lessons: GU-C152-chiththi, GU-C152-maramat, GU-C152-bhadun, GU-C152-malik-chiththi, GU-R152-first-pass-future-notices, GU-R152-second-pass-core

\chapter{Reading a note about a repair}
\label{ch:gu-reading-a-note-about-a-repair}
\hlchaptermodality{152}

\begin{chapteropening}
\textbf{By the end of this chapter:} \emph{I can read a short note about a repair and the rent, find what it asks of me, and answer it.}

Read a 52-word note from a landlord about a repair and the rent, find the one question it asks, and plan the reply.
\end{chapteropening}

\section[\texorpdfstring{\gu{ચિઠ્ઠી} (ciṭṭhī)}{ચિઠ્ઠી (ciṭṭhī)}]{\gu{ચિઠ્ઠી} (ciṭṭhī) — a note}

\label{lesson:GU-C152-chiththi}



\begin{quote}
Before the new one: say \emph{a posted notice}, then \emph{the shop is open only in the morning}.
\end{quote}

\subsection*{You'll want to know: \gu{ચિઠ્ઠી}}
\textbf{\gu{ચિઠ્ઠી}} (\emph{ciṭṭhī}) — \textquotedblleft{}a note, a short written message\textquotedblright{}.

A notice is for everyone; a \textbf{\gu{ચિઠ્ઠી}} is for one person: a line left under a door, a message from a neighbour.

\begin{quote}
\textbf{\gu{મેં} \gu{તેને} \gu{ચિઠ્ઠી} \gu{લખી}.} — \emph{mẽ tene ciṭṭhī lakhī.} — I wrote her a note. \textbf{\gu{ચિઠ્ઠી} \gu{ક્યાં} \gu{છે}?} — \emph{ciṭṭhī kyā̃ che?} — Where is the note?
\end{quote}

\textbf{\gu{લખી}} ends in \textbf{-\gu{ી}} because \textbf{\gu{ચિઠ્ઠી}} is feminine. The \textbf{\gu{ઠ્ઠ}} in the middle is \textbf{\gu{ઠ}} twice, joined by a virama: a doubled, breathy, tongue-curled \emph{ṭṭh}.

\begin{cousinweb}
\textbf{\gu{ચિઠ્ઠી}} (\emph{ciṭṭhī}) is inherited. CDIAL reconstructs an Indo-Aryan \emph{ciṣṭa}, \textquotedblleft{}message\textquotedblright{}, and lists Gujarati \emph{ciṭṭhī}, \textquotedblleft{}message\textquotedblright{}, under it, beside Hindi \emph{ciṭṭhī}, \textquotedblleft{}letter\textquotedblright{} (CDIAL 4832).
\end{cousinweb}

\begin{grammarlens}[title={What you've built}]
One more word for the notes people leave for each other.
\end{grammarlens}

\subsection*{Your turn}
Say these aloud:

\begin{itemize}
\raggedright
  \item \emph{ciṭṭhī}
  \item \emph{mẽ tene ciṭṭhī lakhī}
  \item \emph{ciṭṭhī kyā̃ che?}
\end{itemize}

\begin{tcolorbox}[breakable,colback=teal!4,colframe=teal!35!black,title={Before you move on}]
What does \textbf{\gu{ચિઠ્ઠી}} mean? (\emph{A note}.) Who is it for, compared with a notice? (One person, not everyone.)
\end{tcolorbox}
\section[\texorpdfstring{\gu{મરામત} (marāmat)}{મરામત (marāmat)}]{\gu{મરામત} (marāmat) — a repair}

\label{lesson:GU-C152-maramat}



\begin{quote}
Before the new one: say \emph{I wrote her a note}, then ask \emph{where is the note?}
\end{quote}

\subsection*{You'll want to know: \gu{મરામત}}
\textbf{\gu{મરામત}} (\emph{marāmat}) — \textquotedblleft{}a repair, repair work\textquotedblright{}.

It is the job of putting something right: a window, a door, a road.

\begin{quote}
\textbf{\gu{આવતીકાલે} \gu{મરામત} \gu{થશે}.} — \emph{āvtīkāle marāmat thaśe.} — The repair will be done tomorrow. \textbf{\gu{તે} \gu{મરામત} \gu{કરશે}.} — \emph{te marāmat karśe.} — He will do the repair.
\end{quote}

\textbf{\gu{થશે}} and \textbf{\gu{કરશે}} split the work the way English \emph{be done} and \emph{do} do: \textbf{\gu{થશે}} when the repair happens, \textbf{\gu{કરશે}} when someone does it.

\begin{cousinweb}
\textbf{\gu{મરામત}} (\emph{marāmat}) came into Gujarati through Persian from Arabic \emph{maramma}, \textquotedblleft{}a repairing\textquotedblright{}. Its root is \emph{r-m-m}, \textquotedblleft{}to repair, to put into a sound state\textquotedblright{}, said in Lane of a worn rope, a house or a wall (Lane).
\end{cousinweb}

\begin{grammarlens}[title={What you've built}]
One more word for the notes people leave for each other.
\end{grammarlens}

\subsection*{Your turn}
Say these aloud:

\begin{itemize}
\raggedright
  \item \emph{marāmat}
  \item \emph{āvtīkāle marāmat thaśe}
  \item \emph{te marāmat karśe}, he will do the repair
\end{itemize}

\begin{tcolorbox}[breakable,colback=teal!4,colframe=teal!35!black,title={Before you move on}]
What does \textbf{\gu{મરામત}} mean? (\emph{A repair}.) Which verb says someone does it? (\textbf{\gu{કરશે}}.)
\end{tcolorbox}
\section[\texorpdfstring{\gu{ભાડું} (bhāḍũ)}{ભાડું (bhāḍũ)}]{\gu{ભાડું} (bhāḍũ) — rent}

\label{lesson:GU-C152-bhadun}



\begin{quote}
Before the new one: say \emph{the repair will be done tomorrow}, then \emph{he will do the repair}.
\end{quote}

\subsection*{You'll want to know: \gu{ભાડું}}
\textbf{\gu{ભાડું}} (\emph{bhāḍũ}) — \textquotedblleft{}rent, a fare\textquotedblright{}.

It is the money paid for using something that is not yours: a room, a house, a seat on a bus.

\begin{quote}
\textbf{\gu{ભાડું} \gu{કેટલું} \gu{છે}?} — \emph{bhāḍũ keṭlũ che?} — How much is the rent? How much is the fare? \textbf{\gu{મેં} \gu{ભાડું} \gu{આપ્યું}.} — \emph{mẽ bhāḍũ āpyũ.} — I paid the rent.
\end{quote}

\textbf{\gu{ભાડું}} is neuter, so \textbf{\gu{આપ્યું}}.

\begin{cousinweb}
\textbf{\gu{ભાડું}} (\emph{bhāḍũ}) is inherited. CDIAL lists the Gujarati word under a reconstructed \emph{bhārta}, \textquotedblleft{}pertaining to servants\textquotedblright{}, which Middle Indo-Aryan continued as \emph{bhāṭa}, \textquotedblleft{}wages\textquotedblright{} (CDIAL 9468). Pay for a servant's work became pay for a room.
\end{cousinweb}

\begin{grammarlens}[title={What you've built}]
One more word for the notes people leave for each other.
\end{grammarlens}

\subsection*{Your turn}
Say these aloud:

\begin{itemize}
\raggedright
  \item \emph{bhāḍũ}
  \item \emph{bhāḍũ keṭlũ che?}
  \item \emph{mẽ bhāḍũ āpyũ}
\end{itemize}

\begin{tcolorbox}[breakable,colback=teal!4,colframe=teal!35!black,title={Before you move on}]
What does \textbf{\gu{ભાડું}} mean? (\emph{Rent, a fare}.) Why \textbf{\gu{આપ્યું}}? (\textbf{\gu{ભાડું}} is neuter.)
\end{tcolorbox}
\section[(reading)]{(reading) — a note under the door}

\label{lesson:GU-C152-malik-chiththi}



\begin{quote}
Say \emph{how much is the rent?}, then \emph{the repair will be done tomorrow}. And the word for \emph{owner}: \textbf{\gu{માલિક}}.

The shop notice spoke to everyone. This note speaks to one person, Mira, and asks her for something.
\end{quote}

\begin{tcolorbox}[breakable,skin=enhanced,colback=orange!6,colframe=orange!45!black,boxrule=0.5pt,arc=1mm,left=6pt,right=6pt,top=4pt,bottom=4pt,fonttitle=\bfseries,title={Read this}]
\begin{quote}
\textbf{\gu{ચિઠ્ઠી}}
\end{quote}

\begin{quote}
\textbf{\gu{નમસ્તે} \gu{મીરા} \gu{બહેન},} \textbf{\gu{હું} \gu{અરુણ} \gu{છું}, \gu{માલિક}.} \textbf{\gu{ગઈકાલે} \gu{તમે} \gu{ફોન} \gu{કર્યો}: \gu{તમારી} \gu{બારી} \gu{ખુલ્લી} \gu{રહે} \gu{છે}.} \textbf{\gu{આવતીકાલે} \gu{સવારે} \gu{મરામત} \gu{થશે}.} \textbf{\gu{એક} \gu{માણસ} \gu{આવશે}. \gu{તે} \gu{બારી} \gu{અને} \gu{દરવાજો} \gu{જોશે}.} \textbf{\gu{માણસ} \gu{બપોરે} \gu{જશે}.} \textbf{\gu{શું} \gu{તમે} \gu{ઘરે} \gu{હશો}? \gu{હું} \gu{પણ} \gu{આવીશ}.} \textbf{\gu{ભાડું} \gu{પાંચ} \gu{હજાર} \gu{રૂપિયા} \gu{છે}.} \textbf{\gu{સોમવારે} \gu{હું} \gu{ભાડું} \gu{લઈશ}.} \textbf{\gu{ચિંતા} \gu{ન} \gu{કરો}.} \textbf{\gu{આભાર}, \gu{અરુણ}}
\end{quote}

\emph{Read it:} the note once, straight through, without translating

\emph{Read it:} it a second time, and find the one line that is a question
\end{tcolorbox}

\subsection*{What to know first}
\textbf{\gu{ખુલ્લી} \gu{રહે} \gu{છે}} is \emph{stays open}: the window will not shut, and that is the reason for the note. It runs in time order: yesterday's call in the past, \textbf{\gu{ફોન} \gu{કર્યો}}; tomorrow's repair in the future, \textbf{\gu{થશે}}, \textbf{\gu{આવશે}}, \textbf{\gu{જશે}}; and one question in the middle for Mira to answer.

\subsection*{Your turn}
\begin{itemize}
\raggedright
  \item \emph{Answer:} who will come, and when will he leave? (A man; at midday.)
  \item \emph{Answer:} the rent, and the day Arun collects it. (Five thousand rupees; Monday.)
  \item \emph{Answer:} Mira will be at work tomorrow morning. What should her reply say? (No, and why, as a plan: \textbf{\gu{ના}. \gu{હું} \gu{કામ} \gu{કરવાની} \gu{છું}.})
\end{itemize}

\begin{tcolorbox}[breakable,colback=teal!4,colframe=teal!35!black,title={Before you move on}]
Who is the note for? (One person, Mira.) What does it ask her? (Whether she will be at home tomorrow morning.)
\end{tcolorbox}
\section[(dialogue)]{The future, plans and notices — a review}

\label{lesson:GU-R152-first-pass-future-notices}



\begin{quote}
Say what Arun's note asked Mira, and how much the rent was.

\begin{itemize}
\raggedright
  \item \emph{Write it:} \textbf{\gu{ઊ}} once from memory
\end{itemize}
\end{quote}

\subsection*{Your turn: the future and plans}
Say these aloud:

\begin{itemize}
\raggedright
  \item \emph{I will speak Gujarati}; \emph{we will meet}; \emph{she will come}; \emph{when will you come?}
  \item \emph{I will go}; \emph{I will take the medicine}; \emph{it will be late}; \emph{she will be at home}
  \item \emph{I'm sleepy}; \emph{I'm going tomorrow}, as a plan; \emph{what are you going to do?}
  \item \emph{I'm not going to come}; \emph{perhaps she'll come}
\end{itemize}

\subsection*{Your turn: the notice}
\begin{itemize}
\raggedright
  \item \emph{Say it:} \emph{a posted notice}; \emph{the shop is open}; \emph{only in the morning}; \emph{a public holiday}
  \item \emph{Answer:} on a notice, which words do you hunt for first? (The \emph{whens}.)
\end{itemize}

\begin{tcolorbox}[breakable,colback=teal!4,colframe=teal!35!black,title={Before you move on}]
\emph{I will speak}? (\textbf{\gu{હું} \gu{બોલીશ}}.) \emph{When will you come?} (\textbf{\gu{તમે} \gu{ક્યારે} \gu{આવશો}?}) \emph{It will be late}? (\textbf{\gu{મોડું} \gu{થશે}}.) \emph{I'm sleepy}? (\textbf{\gu{મને} \gu{ઊંઘ} \gu{આવે} \gu{છે}}.) \emph{I'm going tomorrow}? (\textbf{\gu{હું} \gu{કાલે} \gu{જવાનો} \gu{છું}}, or \textbf{\gu{જવાની} \gu{છું}}.) \emph{Perhaps}? (\textbf{\gu{કદાચ}}.) \emph{Open}? (\textbf{\gu{ખુલ્લું}}.) \emph{Public}? (\textbf{\gu{જાહેર}}.)
\end{tcolorbox}
\section[(dialogue)]{The whole A2 core — a second pass}

\label{lesson:GU-R152-second-pass-core}



\begin{quote}
Say \emph{a note}, then \emph{a repair}.
\end{quote}

\subsection*{Your turn: the note}
\begin{itemize}
\raggedright
  \item \emph{Say it:} \emph{I wrote her a note}; \emph{the repair will be done tomorrow}; \emph{how much is the rent?}
  \item \emph{Answer:} what did Arun's note ask Mira? (Whether she will be at home tomorrow morning.)
  \item \emph{Answer:} a notice or a note: which one speaks to one person? (A note: \textbf{\gu{ચિઠ્ઠી}}.)
\end{itemize}

\subsection*{Your turn: one sweep}
Say these aloud:

\begin{itemize}
\raggedright
  \item \emph{I don't speak Gujarati}; \emph{nobody is here}; \emph{are you coming or not?}
  \item \emph{their brother}; \emph{I got worried}; \emph{I came late}; \emph{I did the work}
  \item \emph{I drank tea}; \emph{she ate rice}; \emph{he took the medicine}
  \item \emph{they will speak Gujarati}; \emph{when will you come?}; \emph{what are you going to do tomorrow?}
\end{itemize}

\begin{tcolorbox}[breakable,colback=teal!4,colframe=teal!35!black,title={Before you move on}]
\emph{A note}? (\textbf{\gu{ચિઠ્ઠી}}.) \emph{A repair}? (\textbf{\gu{મરામત}}.) \emph{Rent}? (\textbf{\gu{ભાડું}}.) Which tense did Arun use for the repair, and why? (The future, \textbf{\gu{થશે}}: it has not happened yet.)
\end{tcolorbox}
